\documentclass[../labs.tex]{subfiles}
\begin{document}

\chapter{Laboratory 02: Hugging Face ecosystem, encoder-decoder models, and attention mechanisms}
\label{ch:lab02_huggingface_attention}

\section{Laboratory overview and macro perspective in LLM systems}

In Laboratory~01, we established the fundamental computational substrate of modern deep learning: multidimensional tensors, strided memory layouts, automatic differentiation (autograd) via dynamically constructed Directed Acyclic Graphs (DAGs), and the iterative gradient descent optimization loop. While these primitives allow us to construct arbitrary multilayer perceptrons and convolutional backbones, language modeling introduces a fundamentally different challenge: processing variable-length discrete symbolic sequences where the meaning of any individual token depends dynamically on its contextual relationship with every other token in the sequence.

Laboratory~02 bridges this foundational gap by transitioning from atomic tensor algebra to production-grade \textbf{Transformer architectures} and the \textbf{Attention Mechanism}---the foundational innovation introduced by Vaswani et al.~\cite{vaswani2017} that powers all modern Large Language Models.

\subsection{From tensor primitives to attention-driven foundation models}

To appreciate why the attention mechanism revolutionized deep learning, consider the limitations of the computational layers explored in Laboratory~01:
\begin{enumerate}
    \item \textbf{Linear (Dense) Layers}: A standard linear layer $\vy = \mW \vx + \vb$ computes static affine transformations with fixed input dimensionality. It possesses no inherent mechanism to handle variable-length sequences or to model position-dependent interactions without concatenating all tokens into an unwieldy, parameter-heavy vector.
    \item \textbf{Recurrent Neural Networks (RNNs/LSTMs)}: While historically deployed for sequential text, recurrent models process tokens sequentially ($\vh_t = f(\vh_{t-1}, \vx_t)$). This sequential data dependency creates an insurmountable training bottleneck: forward passes and backpropagation through time (BPTT) cannot be parallelized across sequence lengths $T$ on massively parallel hardware accelerators (GPUs/TPUs). Furthermore, recurrent hidden states suffer from information bottlenecks and vanishing/exploding gradients when bridging long-range token dependencies~\cite{hochreiter1997,pascanu2013difficulty}.
    \item \textbf{Attention Mechanisms}: Attention circumvents both limitations by computing direct, pairwise similarity scores between all tokens in a sequence simultaneously. This enables $O(1)$ path length between any two arbitrary positions regardless of sequential distance, while executing entirely as batched, highly optimized matrix multiplications ($\mat{Q} \mat{K}^\top \mat{V}$) perfectly matched to modern tensor cores~\cite{vaswani2017,dao2022flashattention}.
\end{enumerate}

\subsection{Taxonomy of Transformer architectures: Encoder-decoder, causal decoder, and masked encoder}

Not all Transformer models are structured identically. Depending on the sequence-to-sequence transformation required, the literature divides Transformer topologies into three primary structural archetypes, summarized in Table~\ref{tab:transformer_archetypes}.

\begin{table}[htbp]
\centering
\small
\renewcommand{\arraystretch}{1.3}
\begin{tabularx}{\textwidth}{l p{3.2cm} p{3.2cm} X}
\toprule
\textbf{Architecture} & \textbf{Attention Masking} & \textbf{Canonical Models} & \textbf{Target Applications} \\
\midrule
\textbf{Encoder-Only} & Full bidirectional (unmasked) across all sequence tokens & BERT~\cite{devlin2019}, RoBERTa, CodeBERT~\cite{feng2020codebert} & Token classification, masked language modeling, sentence embeddings, code search. \\
\textbf{Decoder-Only} & Causal triangular mask (tokens attend strictly to past positions) & GPT-2/3/4~\cite{radford2019,brown2020}, LLaMA~\cite{touvron2023llama}, Mistral & Autoregressive text generation, open-ended code synthesis, conversational chat. \\
\textbf{Encoder-Decoder} & Bidirectional in encoder, causal in decoder, cross-attention interconnect & T5~\cite{raffel2020t5}, BART, Original Transformer~\cite{vaswani2017} & Text-to-text transformation, neural machine translation, code summarization, refactoring. \\
\bottomrule
\end{tabularx}
\caption{Taxonomy of Transformer architectural paradigms across attention topology and downstream tasks.}
\label{tab:transformer_archetypes}
\end{table}

This laboratory directly investigates the \textbf{Encoder-Decoder paradigm} using Google's T5 (Text-to-Text Transfer Transformer)~\cite{raffel2020t5}, before isolating and building the underlying \textbf{Scaled Dot-Product Attention mechanism} from first principles in pure PyTorch.

\subsection{Pedagogical roadmap and learning outcomes}

Laboratory~02 is divided into two deeply complementary computational modules:
\begin{enumerate}
    \item \textbf{Module 1: Hugging Face Ecosystem and T5 Architecture Exploration (\texttt{text-01-t5.ipynb})}:
    \begin{itemize}
        \item Ingesting and tokenizing text using the Hugging Face \texttt{transformers} library.
        \item Mastering SentencePiece subword tokenization, vocabulary structures, and the whitespace-delimited underscore prefix (\texttt{\_}).
        \item Analyzing sequence batching, zero-padding, attention masking, and GPU memory overheads.
        \item Inspecting T5's internal parameter organization, multi-head dimensions ($768 = 12 \times 64$), and weight-sharing mechanisms across embeddings and linear heads.
        \item Geometric inspection of latent embeddings via Principal Component Analysis (PCA) and Cosine Similarity matrices.
        \item Implementing the manual autoregressive generation loop step-by-step, contrasting it with optimized KV-cached generation (\texttt{model.generate()}).
        \item Extracting, dissecting, and visualizing internal cross-attention tensors across all 12 layers and attention heads.
    \end{itemize}
    \item \textbf{Module 2: Scaled Dot-Product Attention from Scratch on Synthetic Sequences (\texttt{text-02-attention.ipynb})}:
    \begin{itemize}
        \item Formulating discontinuous periodic time series prediction as a sequence regression task.
        \item Constructing custom PyTorch \texttt{Dataset} and \texttt{DataLoader} pipelines with randomized jump discontinuities.
        \item Implementing custom \texttt{ScaledDotProductAttention} and multi-parameter \texttt{AttentionHead} modules from scratch.
        \item Mathematically deriving the $\frac{1}{\sqrt{\dk}}$ scaling factor and its variance-stabilization proof.
        \item Conducting a rigorous empirical benchmark comparing \textbf{Relative Positional Encodings} (Shaw et al.~\cite{shaw2018} / T5 style) against \textbf{Absolute Sinusoidal Positional Encodings} (Vaswani et al.~\cite{vaswani2017}).
        \item Proving why relative positional representations converge to near-zero loss ($0.0363$) on periodic invariants where absolute encodings completely stagnate ($5.4174$).
    \end{itemize}
\end{enumerate}

\begin{figure}[htbp]
\centering
\begin{tikzpicture}[
    scale=0.88,
    every node/.style={transform shape},
    block/.style={rectangle, draw=politoNavy, fill=skyBlue!30, very thick, rounded corners=4pt, minimum height=11mm, minimum width=24mm, font=\small\bfseries, align=center},
    focusblock/.style={rectangle, draw=deepdiveTeal, fill=deepdiveTealBg, very thick, rounded corners=4pt, minimum height=12mm, minimum width=26mm, font=\small\bfseries, align=center},
    arr/.style={-{Stealth[scale=1.1]}, thick, draw=bodyDark}
]

\node[block] (tok) at (0, 0) {1. Raw Text\\Sentence};
\node[focusblock] (sp) at (3.2, 0) {2. SentencePiece\\Subwords (\_)\\{\scriptsize (Module 1)}};
\node[focusblock] (pad) at (6.6, 0) {3. Batch Padding\\\& Attention Mask\\{\scriptsize (Module 1)}};
\node[focusblock] (emb) at (10.0, 0) {4. Shared Embedding\\Space (32128x768)\\{\scriptsize (Module 1)}};
\node[focusblock] (enc) at (13.4, 0) {5. Bidirectional\\T5 Encoder\\{\scriptsize (Module 1)}};

\node[block] (out) at (0, -2.8) {10. Decoded Output\\Text String};
\node[focusblock] (gen) at (3.2, -2.8) {9. Greedy / Beam\\Autoregressive Loop\\{\scriptsize (Module 1)}};
\node[focusblock] (head) at (6.6, -2.8) {8. Tied LM Head\\Logits ($B, T, |\vocab|$)};
\node[focusblock] (cross) at (10.0, -2.8) {7. Cross-Attention\\Alignments\\{\scriptsize (Module 1)}};
\node[focusblock] (dec) at (13.4, -2.8) {6. Causal T5\\Decoder Block\\{\scriptsize (Module 1)}};

\draw[arr] (tok) -- (sp);
\draw[arr] (sp) -- (pad);
\draw[arr] (pad) -- (emb);
\draw[arr] (emb) -- (enc);
\draw[arr] (enc) -- (dec);
\draw[arr] (dec) -- (cross);
\draw[arr] (cross) -- (head);
\draw[arr] (head) -- (gen);
\draw[arr] (gen) -- (out);

\end{tikzpicture}
\caption{End-to-end dataflow in the Hugging Face T5 Sequence-to-Sequence pipeline explored throughout Laboratory~02.}
\label{fig:lab02_t5_pipeline}
\end{figure}

% ==============================================================================
% SECTION 2: TOKENIZATION PIPELINES AND BATCH PADDING
% ==============================================================================
\section{Tokenization pipelines, subwords, and batch padding mechanics}
\label{sec:lab02_tokenization}

Large Language Models do not process raw ASCII or UTF-8 text strings directly. Instead, text must undergo deterministic conversion into a discrete sequence of integer indices drawn from a predefined vocabulary $\vocab$, a process termed \textbf{tokenization}.

\subsection{The role of tokenizers in language models}

In early Natural Language Processing pipelines, tokenization relied on whitespace splitting or heuristic rule-based stemmers. These approaches faced two catastrophic failure modes:
\begin{enumerate}
    \item \textbf{Out-Of-Vocabulary (OOV) Explosion}: Natural languages exhibit virtually unbounded vocabularies due to morphological compounding, slang, spelling variations, and programming language identifiers (\texttt{getUserById}, \texttt{is\_valid}). Word-level vocabularies quickly expand into millions of tokens while still failing on unseen words.
    \item \textbf{Character-Level Dilution}: While character-level models eliminate OOV tokens, they drastically inflate sequence lengths ($T$), causing self-attention computational complexity ($O(T^2)$) and memory consumption to explode, while diluting semantic signal across individual characters.
\end{enumerate}

Modern foundational models resolve this dilemma via \textbf{Subword Tokenization algorithms} (such as Byte-Pair Encoding (BPE)~\cite{sennrich2016}, WordPiece, and SentencePiece Unigram~\cite{kudo2018}). Frequent words remain single atomic tokens, while rare, compound, or domain-specific words are decomposed into compact subword fragments.

\subsection{SentencePiece and the whitespace-delimited underscore prefix}

The T5 model utilizes Google's \textbf{SentencePiece} tokenizer library with a Unigram language model vocabulary of $32{,}128$ tokens~\cite{kudo2018,raffel2020t5}. Unlike traditional tokenizers that treat whitespace as an external delimiter discarded before tokenization, SentencePiece treats whitespace as a first-class character, represented by a special meta-symbol: the unicode underscore prefix \texttt{\_} (specifically U+2581 \texttt{\_}).

\begin{intuition}[The Underscore Prefix as a Reversible Word Boundary]
Why does SentencePiece prepend an underscore \texttt{\_} to tokens instead of storing separate space characters?

In standard tokenization, converting \texttt{"hello world"} into tokens and then decoding might lose exact whitespace information (e.g., whether there was a single space, multiple spaces, tabs, or no space before punctuation). By explicitly prepending \texttt{\_} to any subword that was preceded by whitespace, tokenization becomes strictly \textbf{lossless and reversible} (\textit{detokenization} requires only concatenating tokens and replacing \texttt{\_} with a space). 

Crucially, as emphasized during the course lecture by teaching assistant Claudio Savelli, the underscore also alters the token's semantic identity:
\begin{itemize}
    \item \texttt{"\_pre"}: Represents the token \texttt{pre} appearing at the beginning of a word or after whitespace (e.g., \texttt{"\_pre-requisite"}).
    \item \texttt{"pre"}: Represents the identical character sequence occurring inside a word (e.g., within \texttt{"com\_pre\_hend"}).
\end{itemize}
The neural network assigns entirely distinct embedding vectors to \texttt{"\_pre"} and \texttt{"pre"}, allowing it to contextualize morphological prefixes differently from internal syllables.
\end{intuition}

\subsection{Batching, padding, and attention masks}

Deep learning models achieve their computational throughput by evaluating multiple sequences simultaneously in a batch ($B > 1$). Because GPU kernels (such as cuBLAS GEMM) execute on rectangular multidimensional tensors of fixed shape $[B, T, d_{\text{model}}]$, all sequences in a given batch must possess identical length $T$.

When batching sentences of varying natural lengths, two operations are required:
\begin{enumerate}
    \item \textbf{Zero-Padding (\texttt{padding=True})}: The longest sequence in the batch defines the maximum length $T_{\text{max}}$. All shorter sequences are padded with a designated pad token (\texttt{pad\_token\_id = 0} in T5) until they reach $T_{\text{max}}$.
    \item \textbf{Attention Masking}: Adding dummy pad tokens introduces a severe correctness hazard: self-attention would ordinarily compute query-key dot products between legitimate semantic tokens and meaningless pad tokens, polluting representations and skewing softmax distributions. To prevent this, the tokenizer produces a binary \textbf{Attention Mask} tensor $\mM \in \{0, 1\}^{B \times T}$:
    \begin{equation}
        M_{b, t} = \begin{cases}
            1, & \text{if position } t \text{ is an authentic input token}, \\
            0, & \text{if position } t \text{ is a padded token}.
        \end{cases}
    \end{equation}
    In the attention mechanism, positions where $M_{b, t} = 0$ have $-\infty$ added to their pre-softmax attention logits, ensuring that $\softmax(-\infty) = 0$ and completely nullifying any information flow to or from padded positions.
\end{enumerate}

\begin{examinsight}[Memory Saturation and Dynamic Bucketing]
Padding introduces a critical GPU memory and computational inefficiency. If a batch contains 31 short sentences of length 10 and one outlier sentence of length 512, all 31 short sequences are padded to length 512! Over $98\%$ of the allocated tensor elements represent non-informative zeros, saturating GPU VRAM and performing wasteful $O(T_{\text{max}}^2)$ attention calculations.

In production LLM training pipelines (and explored further in Laboratory~03), engineers counteract this via:
\begin{enumerate}
    \item \textbf{Truncation (\texttt{truncation=True, max\_length=K})}: Hard-capping sequence length to prevent single outliers from blowing up batch shapes.
    \item \textbf{Dynamic Batching / Length Bucketing}: Sorting dataset sequences by token length prior to batch construction, ensuring that sequences within each batch possess nearly identical lengths, minimizing zero-padding overhead.
\end{enumerate}
\end{examinsight}

\subsection{Verified code: Tokenization, padding, and batch decoding}

Listing~\ref{lst:lab02_tokenizer_init} displays the verified implementation from Cell~2 of \texttt{text-01-t5.ipynb}, loading the pretrained T5 tokenizer and inspecting its class hierarchy.

\begin{lstlisting}[language=Python, caption={Loading T5 tokenizer and verifying class type (Cells 2--5).}, label={lst:lab02_tokenizer_init}]
from transformers import AutoTokenizer

model_name = "google-t5/t5-base"

# Instantiate tokenizer from Hugging Face Hub
tokenizer = AutoTokenizer.from_pretrained(model_name)

# Verified check: print the type of the tokenizer
print(type(tokenizer))
# Output: <class 'transformers.models.t5.tokenization_t5_fast.T5TokenizerFast'>

# Tokenize a basic sentence
sentence = "hello, this is a sentence!"
tokens = tokenizer(sentence)
print(tokens)
\end{lstlisting}

Executing this cell produces the token dictionary:
\begin{lstlisting}[style=terminalOutput]
{'input_ids': [21820, 6, 44, 19, 3, 9, 7142, 55, 1], 'attention_mask': [1, 1, 1, 1, 1, 1, 1, 1, 1]}
\end{lstlisting}

Notice the trailing token ID \texttt{1}. In T5, token ID \texttt{1} represents the End-Of-Sequence sentinel (\texttt{</s>}), explicitly appended by the tokenizer to mark sequence completion.

We can inspect the vocabulary size and probe the internal token-to-ID mapping using Cells~10--28 of the notebook:

\begin{lstlisting}[language=Python, caption={Probing vocabulary structure and special token IDs (Cells 10--28).}, label={lst:lab02_vocab_probe}]
import random 

vocabulary = tokenizer.get_vocab()
print("Total vocabulary size:", len(vocabulary)) # Output: 32128

# Reverse mapping: token_id -> string
reverse_vocab = {v: k for k, v in vocabulary.items()}

# Inspect special tokens
print("Special tokens:", tokenizer.eos_token, tokenizer.pad_token, tokenizer.bos_token)
# Output: </s> <pad> None

print("Special token IDs:", tokenizer.eos_token_id, tokenizer.pad_token_id, tokenizer.bos_token_id)
# Output: 1 0 None

# Inspect subword tokenization with whitespace variations
print(tokenizer.tokenize("hello, this is a sentence!"))
# Output: [' hello', ',', ' this', ' is', ' ', 'a', ' sentence', '!']

print(tokenizer.tokenize("hello    ,world"))
# Output: [' hello', ' ', ',', 'world']
print(tokenizer.tokenize("hello    , world"))
# Output: [' hello', ' ', ',', ' world']
\end{lstlisting}

In the lecture walkthrough, TA Claudio Savelli highlighted that SentencePiece collapses multiple consecutive spaces into a single \texttt{\_} meta-token followed by standard subwords, demonstrating the whitespace-normalizing properties of the Unigram tokenizer.

Next, Listing~\ref{lst:lab02_batch_padding} demonstrates the verified solution for batch tokenization with automatic padding, contrasting standard list outputs with PyTorch tensor outputs.

\begin{lstlisting}[language=Python, caption={Batch tokenization with zero-padding and tensor formatting (Cells 30--36).}, label={lst:lab02_batch_padding}]
sentences = [
    "this is the first sentence",
    "instead, this is the second sequence!"
]

# Verified Solution: Tokenize sentences with padding as PyTorch tensors
tokens = tokenizer(sentences, padding=True, return_tensors="pt")
print("Input IDs Tensor:")
print(tokens["input_ids"])
print("Attention Mask Tensor:")
print(tokens["attention_mask"])

# Decode the batch to inspect padding tokens
decoded_batch = tokenizer.batch_decode(tokens["input_ids"])
print("Decoded batch:", decoded_batch)
\end{lstlisting}

The execution output demonstrates how SentencePiece pads the shorter first sentence to match the nine-token length of the second sentence:
\begin{lstlisting}[style=terminalOutput]
Input IDs Tensor:
tensor([[  48,   19,    8,  166, 7142,    1,    0,    0,    0],
        [1039,    6,   48,   19,    8,  542, 3817,   55,    1]])

Attention Mask Tensor:
tensor([[1, 1, 1, 1, 1, 1, 0, 0, 0],
        [1, 1, 1, 1, 1, 1, 1, 1, 1]])

Decoded batch: ['this is the first sentence</s><pad><pad><pad>', 'instead, this is the second sequence!</s>']
\end{lstlisting}

Notice that while \texttt{tokens["input\_ids"]} contains trailing zeros (\texttt{<pad>}), \texttt{tokens["attention\_mask"]} sets exactly those positions to \texttt{0}, protecting downstream attention layers from calculating spurious correlations.

% ==============================================================================
% SECTION 3: ANATOMY OF THE T5 ARCHITECTURE AND EMBEDDINGS
% ==============================================================================
\section{Anatomy of the T5 architecture and embedding geometry}
\label{sec:lab02_t5_anatomy}

Having parsed discrete text into clean tensor batches, we load the full T5 foundation model: \texttt{google-t5/t5-base}.

\subsection{The T5 unified text-to-text paradigm}

Introduced by Raffel et al.~\cite{raffel2020t5}, T5 reconceptualizes all natural language processing tasks into a unified \textbf{Text-to-Text} formulation. Whether performing machine translation, document summarization, grammatical error correction, or sentiment classification, both the input and the target output are represented as text strings:
\begin{itemize}
    \item \textbf{Machine Translation}: \texttt{"translate English to German: That is good."} $\longrightarrow$ \texttt{"Das ist gut."}
    \item \textbf{Summarization}: \texttt{"summarize: State of the union address..."} $\longrightarrow$ \texttt{"The President discussed..."}
    \item \textbf{Classification}: \texttt{"sst2 sentence: The movie was gripping."} $\longrightarrow$ \texttt{"positive"}
\end{itemize}
Because the task specification is embedded directly within the input prompt as a natural language prefix, a single unified architecture and identical cross-entropy loss function can be applied across all downstream domains without requiring task-specific classification heads.

\subsection{Model configuration and hyperparameter inspection}

Using Hugging Face's \texttt{AutoModelForSeq2SeqLM}, we instantiate the model and inspect its architecture:

\begin{lstlisting}[language=Python, caption={Loading T5-base model and inspecting architectural configuration (Cells 39--44).}, label={lst:lab02_model_config}]
from transformers import AutoModelForSeq2SeqLM

# Load pretrained sequence-to-sequence model
model = AutoModelForSeq2SeqLM.from_pretrained(model_name)
print(type(model))
# Output: <class 'transformers.models.t5.modeling_t5.T5ForConditionalGeneration'>

# Inspect hyperparameter configuration
cfg = model.config
print(f"Hidden size (d_model): {cfg.d_model}")          # 768
print(f"Number of encoder layers: {cfg.num_layers}")    # 12
print(f"Number of decoder layers: {cfg.num_decoder_layers}") # 12
print(f"Number of attention heads: {cfg.num_heads}")    # 12
print(f"Head dimension (d_kv): {cfg.d_kv}")            # 64
print(f"Feed-forward dimension (d_ff): {cfg.d_ff}")    # 3072
print(f"Vocabulary size: {cfg.vocab_size}")            # 32128
\end{lstlisting}

\begin{deepdive}[Dimensional Multi-Head Alignment in T5]
In T5-base, each hidden token representation possesses dimension $\dmodel = 768$. When entering an attention layer, the token vector is projected into $H = 12$ distinct attention heads. 

Each head operates on a projected subspace of dimensionality $\dk = 64$. Notice the exact algebraic equality:
\begin{equation}
    H \times \dk = 12 \times 64 = 768 = \dmodel
\end{equation}
This design guarantees that projecting tokens into 12 distinct representational subspaces and concatenating the multi-head outputs produces a vector of dimension exactly matching $\dmodel$, allowing residual skip connections ($\vx + \text{Sublayer}(\vx)$) to be added without projection matrices.
\end{deepdive}

\subsection{Weight sharing across encoder, decoder, and language modeling head}

Inspecting the instantiated model reveals a critical memory optimization: \textbf{Weight Tying (Weight Sharing)}. In T5, the embedding layer \texttt{model.shared} provides the token embedding lookup table for both the encoder and the decoder. Furthermore, the final output linear projection layer (\texttt{model.lm\_head}), which projects decoder hidden states back to vocabulary logits of dimension $32{,}128$, reuses the exact same weight matrix!

Listing~\ref{lst:lab02_weight_sharing} displays Cell~49 of \texttt{text-01-t5.ipynb}, verifying this memory sharing using Python pointer identities:

\begin{lstlisting}[language=Python, caption={Verifying weight sharing across encoder, decoder, and LM head (Cells 47--49).}, label={lst:lab02_weight_sharing}]
# Verify memory sharing of embedding layers
is_shared = (id(model.shared) == id(model.encoder.embed_tokens) and 
             id(model.shared) == id(model.decoder.embed_tokens))
print("Are encoder and decoder embedding layers physically shared?", is_shared)
# Output: True

# Verify weight sharing with the LM classification head
print("Weight tying check:", id(model.shared.weight) == id(model.lm_head.weight))
# Output: True
\end{lstlisting}

\begin{examinsight}[Parameter Count Savings via Weight Tying]
Why do modern models tie input embedding weights to output projection weights?
\begin{enumerate}
    \item \textbf{Parameter Economy}: An embedding table of size $|\vocab| \times \dmodel = 32{,}128 \times 768$ contains approximately $24.67$ million parameters. If the encoder embeddings, decoder embeddings, and language modeling head each maintained separate parameters, the model would require $3 \times 24.67 \approx 74$ million parameters solely for token mapping! Weight tying collapses this overhead to a single shared $24.67$M matrix, reducing total model parameter footprint by nearly $50$ million parameters.
    \item \textbf{Representational Duality}: Mathematically, if token $w_i$ has embedding vector $\ve_i$, predicting next-token probability via the dot product $\ve_i^\top \vh_{\text{dec}}$ aligns the geometry of the input semantic space with the output prediction space.
\end{enumerate}
\end{examinsight}

\subsection{Latent embedding space geometry: PCA and cosine similarity}

In Cells~51--54 of \texttt{text-01-t5.ipynb}, the laboratory explores whether the raw un-contextualized vectors stored in \texttt{model.shared} exhibit semantic clustering. We extract embeddings for two semantic clusters: domestic furniture (\texttt{chair}, \texttt{table}, \texttt{plate}, \texttt{knife}, \texttt{spoon}) and domestic animals (\texttt{horse}, \texttt{goat}, \texttt{sheep}, \texttt{cat}, \texttt{dog}).

\begin{lstlisting}[language=Python, caption={Extracting word embeddings and computing PCA and Cosine Similarity (Cells 51--54).}, label={lst:lab02_pca_similarity}]
words = [
    "chair", "table", "plate", "knife", "spoon",
    "horse", "goat", "sheep", "cat", "dog",
]

# Extract token IDs and lookup embeddings in model.shared
word_tokens = tokenizer(words, return_tensors="pt", padding=True)["input_ids"][:, 0]
token_embeddings = model.shared(word_tokens).cpu().detach().numpy()
print("Embedding matrix shape:", token_embeddings.shape) # (10, 768)

# 1. 2D Principal Component Analysis (PCA) projection
from sklearn.decomposition import PCA
import matplotlib.pyplot as plt

pca = PCA(n_components=2)
X_pca = pca.fit_transform(token_embeddings)

fig, ax = plt.subplots(figsize=(5, 5))
ax.scatter(X_pca[:, 0], X_pca[:, 1])
for i, txt in enumerate(words):
    ax.annotate(txt, (X_pca[i, 0], X_pca[i, 1]))
plt.title("PCA Projection of T5 Word Embeddings")
plt.savefig("pca_embeddings.png")

# 2. Pairwise Cosine Similarity Matrix
from sklearn.metrics.pairwise import cosine_similarity

sim_matrix = cosine_similarity(token_embeddings)
fig, ax = plt.subplots(figsize=(5, 5))
cax = ax.imshow(sim_matrix, cmap="coolwarm")
ax.set_xticks(range(len(words)))
ax.set_yticks(range(len(words)))
ax.set_xticklabels(words, rotation=45)
ax.set_yticklabels(words)
fig.colorbar(cax)
plt.title("Cosine Similarity Heatmap")
plt.savefig("cosine_similarity_matrix.png")
\end{lstlisting}

\begin{figure}[htbp]
\centering
\begin{subfigure}[b]{0.48\textwidth}
    \centering
    \includegraphics[width=\textwidth]{figures/lab02/pca_embeddings.png}
    \caption{2D PCA projection of T5 embeddings.}
    \label{fig:t5_pca}
\end{subfigure}
\hfill
\begin{subfigure}[b]{0.48\textwidth}
    \centering
    \includegraphics[width=\textwidth]{figures/lab02/cosine_similarity_matrix.png}
    \caption{Pairwise Cosine Similarity matrix.}
    \label{fig:t5_cosine}
\end{subfigure}
\caption{Geometric structure of T5-base static embedding representations extracted in Laboratory~02. The latent space clearly segregates household items from mammalian animals, confirming semantic clustering prior to contextual attention processing.}
\label{fig:t5_embeddings_geo}
\end{figure}

The resulting visualizations in Figure~\ref{fig:t5_embeddings_geo} demonstrate strong semantic clustering: animals form a tight cluster with high pairwise cosine similarities ($>0.6$), distinctly segregated from household objects.

% ==============================================================================
% SECTION 4: TRANSFORMER ENCODER AND DECODER BLOCKS IN T5
% ==============================================================================
\section{Transformer encoder and decoder blocks in T5}
\label{sec:lab02_blocks}

The complete T5-base model consists of 12 stacked \textbf{Encoder Blocks} and 12 stacked \textbf{Decoder Blocks}.

\subsection{The T5 encoder: Bidirectional self-attention blocks}

Each encoder block (\texttt{model.encoder.block[i]}) contains two sub-layers:
\begin{enumerate}
    \item \textbf{Bidirectional Self-Attention (\texttt{T5LayerSelfAttention})}: Tokens attend to all other tokens in the input sequence simultaneously. Because the entire input prompt is available upfront, no causal masking is applied. The pre-attention normalization follows a pre-LN design without additive bias parameters (T5 uses RMS-like scaling instead of standard LayerNorm~\cite{ba2016layer}).
    \item \textbf{Position-wise Feed-Forward Network (\texttt{T5LayerFF})}: A two-layer MLP projecting representations from $\dmodel = 768 \to \dff = 3072$ using a gated ReLU activation, followed by projection back to $768$.
\end{enumerate}

\subsection{The T5 decoder: Causal self-attention and cross-attention}

Unlike the encoder, each decoder block (\texttt{model.decoder.block[i]}) contains \textbf{three} sub-layers:
\begin{enumerate}
    \item \textbf{Causal Self-Attention (\texttt{T5LayerSelfAttention})}: During autoregressive token generation, the model cannot observe future tokens. A lower-triangular causal mask enforces that position $i$ attends only to positions $j \le i$.
    \item \textbf{Cross-Attention (\texttt{T5LayerCrossAttention})}: This is the critical connective tissue of the encoder-decoder architecture. The Query vectors ($\mQ$) are projected from the decoder's intermediate representations, while the Key ($\mK$) and Value ($\mV$) vectors are projected from the final output representations of the encoder:
    \begin{equation}
        \mQ = \mH_{\text{dec}} \mW_q, \quad \mK = \mH_{\text{enc}} \mW_k, \quad \mV = \mH_{\text{enc}} \mW_v
    \end{equation}
    This allows every generated output token to dynamically inspect and extract information from the input source prompt.
    \item \textbf{Position-wise Feed-Forward Network (\texttt{T5LayerFF})}: Identical in dimension to the encoder MLP ($\dmodel \to 3072 \to \dmodel$).
\end{enumerate}

\begin{figure}[htbp]
\centering
\begin{tikzpicture}[
    scale=0.92,
    every node/.style={transform shape},
    box/.style={rectangle, draw=politoNavy, fill=skyBlue!20, very thick, rounded corners=3pt, minimum width=36mm, minimum height=8mm, font=\small\bfseries, align=center},
    encbox/.style={rectangle, draw=deepdiveTeal, fill=deepdiveTealBg, very thick, rounded corners=3pt, minimum width=36mm, minimum height=8mm, font=\small\bfseries, align=center},
    decbox/.style={rectangle, draw=crimsonRed, fill=examAmberBg, very thick, rounded corners=3pt, minimum width=36mm, minimum height=8mm, font=\small\bfseries, align=center},
    sharedbox/.style={rectangle, draw=intuitionPurple, fill=intuitionPurpleBg, very thick, rounded corners=3pt, minimum width=38mm, minimum height=9mm, font=\small\bfseries, align=center},
    arr/.style={-{Stealth[scale=1.1]}, thick, draw=bodyDark},
    dasharr/.style={-{Stealth[scale=1.1]}, thick, dashed, draw=crimsonRed}
]

% Left: Encoder
\node[sharedbox] (in_emb) at (0, 0) {Shared Token Embeddings\\($32{,}128 \times 768$)};
\node[encbox] (enc_attn) at (0, 1.8) {Bidirectional Self-Attention\\($H=12, \dk=64$)};
\node[encbox] (enc_ffn) at (0, 3.2) {Feed-Forward MLP\\($768 \to 3072 \to 768$)};
\node[draw=deepdiveTealLine, thick, dotted, fit={(enc_attn) (enc_ffn)}, inner sep=4pt, label=left:\textbf{Encoder Block ($\times 12$)}] (enc_block) {};
\node[box] (enc_out) at (0, 4.8) {Encoder Final Memory\\$\mH_{\text{enc}} \in \R^{B \times T_{\text{enc}} \times 768}$};

% Right: Decoder
\node[sharedbox] (dec_emb) at (6.5, 0) {Shared Token Embeddings\\($32{,}128 \times 768$)};
\node[decbox] (dec_causal) at (6.5, 1.4) {Causal Self-Attention\\(Triangular Masked)};
\node[decbox] (dec_cross) at (6.5, 2.7) {Cross-Attention\\($\mQ$ from dec, $\mK,\mV$ from enc)};
\node[decbox] (dec_ffn) at (6.5, 4.0) {Feed-Forward MLP\\($768 \to 3072 \to 768$)};
\node[draw=crimsonRed, thick, dotted, fit={(dec_causal) (dec_cross) (dec_ffn)}, inner sep=4pt, label=right:\textbf{Decoder Block ($\times 12$)}] (dec_block) {};
\node[sharedbox] (lm_head) at (6.5, 5.4) {Tied LM Head Projection\\($768 \to 32{,}128$)};
\node[box] (logits) at (6.5, 6.7) {Vocabulary Logits\\$\in \R^{B \times T_{\text{dec}} \times 32{,}128}$};

% Inputs and Connections
\node[below=0.4cm of in_emb, font=\small] (in_txt) {Source Tokens ($T_{\text{enc}}$)};
\node[below=0.4cm of dec_emb, font=\small] (dec_txt) {Target Tokens ($T_{\text{dec}}$)};

\draw[arr] (in_txt) -- (in_emb);
\draw[arr] (in_emb) -- (enc_attn);
\draw[arr] (enc_attn) -- (enc_ffn);
\draw[arr] (enc_ffn) -- (enc_out);

\draw[arr] (dec_txt) -- (dec_emb);
\draw[arr] (dec_emb) -- (dec_causal);
\draw[arr] (dec_causal) -- (dec_cross);
\draw[arr] (dec_cross) -- (dec_ffn);
\draw[arr] (dec_ffn) -- (lm_head);
\draw[arr] (lm_head) -- (logits);

% Cross-attention link
\draw[dasharr] (enc_out.east) to[out=0, in=180] node[above, font=\scriptsize\bfseries, text=crimsonRed] {Keys \& Values} (dec_cross.west);

% Shared link
\draw[<->, thick, draw=intuitionPurple, dashed] (in_emb.east) -- node[above, font=\scriptsize\bfseries, text=intuitionPurple] {Tied Weights} (dec_emb.west);

\end{tikzpicture}
\caption{The T5 Encoder-Decoder architecture explored in Laboratory~02. The encoder executes bidirectional self-attention across the input prompt; the decoder applies causal self-attention across previously generated tokens and cross-attention over the encoder's final representations; embeddings and classification head are tied.}
\label{fig:t5_full_arch}
\end{figure}

Figure~\ref{fig:t5_full_arch} visually synthesizes the structural dataflow connecting input text, token embeddings, bidirectional encoder layers, cross-attention interconnects, causal decoder layers, and the vocabulary logits.

% ==============================================================================
% SECTION 5: AUTOREGRESSIVE TOKEN GENERATION AND DECODING DYNAMICS
% ==============================================================================
\section{Autoregressive token generation and decoding dynamics}
\label{sec:lab02_generation}

Foundation models generate text \textbf{autoregressively}: token generation proceeds sequentially, where each freshly sampled token is appended to the decoder's context before predicting the subsequent token:
\begin{equation}
    P(y_1, y_2, \dots, y_M \mid x) = \prod_{t=1}^M P(y_t \mid y_{<t}, x)
\end{equation}

\subsection{Logits, vocabularies, and next-token prediction}

When passing inputs through the model, the forward call returns a \texttt{logits} tensor of shape $[B, T_{\text{dec}}, |\vocab|]$. Each element $z_{b, t, v}$ represents the unnormalized log-probability of token ID $v$ at step $t$.

To initialize the generation in T5, the decoder must be primed with an initial start token. In T5, the generation prefix token is the pad token (\texttt{decoder\_input\_ids = [[0]]}).

\subsection{Step-by-step manual autoregressive generation loop}

In Cells~69--85 of \texttt{text-01-t5.ipynb}, the laboratory manually executes the autoregressive loop to translate \texttt{"translate english to german: hello, how are you?"} into German.

\begin{lstlisting}[language=Python, caption={Step-by-step manual autoregressive decoding implementation (Cells 69--85).}, label={lst:lab02_manual_generation}]
import torch

input_sentence = "translate english to german: hello, how are you?"
tokens = tokenizer(input_sentence, return_tensors="pt")

# Step 1: Initialize decoder input with pad token
decoder_input_ids = torch.tensor([[ tokenizer.pad_token_id ]])
print("Initial decoder input:", decoder_input_ids) # tensor([[0]])

# Forward pass through encoder and first decoder step
output = model(**tokens, decoder_input_ids=decoder_input_ids)
print("Logits shape:", output.logits.shape) # torch.Size([1, 1, 32128])

# Extract predicted token at position 0 via greedy argmax
max_proba_token = output.logits[0, 0].argmax()
print(f"Token ID: {max_proba_token.item()} -> Word: '{tokenizer.decode(max_proba_token)}'")
# Output: Token ID: 21820 -> Word: 'Hallo'

# Step 2: Append newly generated token to decoder context
decoder_input_ids = torch.hstack([decoder_input_ids, max_proba_token.view(1, 1)])
print("Updated decoder inputs:", decoder_input_ids) # tensor([[0, 21820]])

# Step 3: Run full manual loop until </s> or max_length
decoder_input_ids = torch.tensor([[ tokenizer.pad_token_id ]])
max_length = 50
i = 0

while i < max_length and decoder_input_ids[0, -1] != tokenizer.eos_token_id:
    output = model(**tokens, decoder_input_ids=decoder_input_ids)
    next_token = output.logits[0, -1].argmax().view(1, 1)
    decoder_input_ids = torch.hstack([decoder_input_ids, next_token])
    i += 1

translated_text = tokenizer.decode(decoder_input_ids[0], skip_special_tokens=True)
print("Manually generated translation:", translated_text)
# Output: "Hallo, wie geht es Ihnen?"
\end{lstlisting}

\subsection{The \texttt{model.generate()} API and Key-Value (KV) caching}

While the manual \texttt{while} loop in Listing~\ref{lst:lab02_manual_generation} illustrates the mathematical logic of autoregression, it is computationally naive. Notice what occurs at iteration $t$:
\begin{enumerate}
    \item The model re-encodes the input sequence (or passes the full encoder states).
    \item More critically, the decoder recomputes self-attention projections ($\mQ, \mK, \mV$) for all previous $t-1$ tokens, even though their causal past has not changed!
\end{enumerate}
Consequently, generating a sequence of length $T_{\text{dec}}$ naively requires $O(T_{\text{dec}}^2)$ attention operations.

Hugging Face's optimized \texttt{model.generate()} method circumvents this recomputation by activating \textbf{Key-Value Caching (KV-Cache)} via the internal argument \texttt{use\_cache=True}:

\begin{lstlisting}[language=Python, caption={High-level optimized generation using \texttt{model.generate()} (Cells 87--88).}, label={lst:lab02_hf_generate}]
# Production generation with KV-cache and beam search / greedy decoding
out_tokens = model.generate(**tokens)
translation = tokenizer.batch_decode(out_tokens, skip_special_tokens=True)
print("Optimized generation:", translation[0])
# Output: "Hallo, wie geht es Ihnen?"
\end{lstlisting}

\begin{examinsight}[KV-Cache Mechanics and Complexity Comparison]
On examinations, students are frequently asked to contrast naive autoregressive generation with KV-cached decoding:
\begin{itemize}
    \item \textbf{Naive Generation (Without Cache)}: At step $t$, the decoder processes a sequence of length $t$. The self-attention matrix requires computing $t \times t$ dot products. Over an entire generation of $M$ tokens, total attention operations scale as:
    \begin{equation}
        \sum_{t=1}^M O(t) = O(M^2)
    \end{equation}
    \item \textbf{Cached Generation (With KV-Cache)}: At step $t$, the decoder passes only the single latest token ($T=1$). Its query $\vq_t$ attends to previously stored key and value vectors ($\mK_{<t}, \mV_{<t}$), and its newly computed $\vk_t, \vv_t$ are appended to the cache. Each step takes $O(t)$ operations for dot products, but avoids re-running linear projections for the previous $t-1$ tokens, drastically slashing wall-clock latency.
\end{itemize}
\end{examinsight}

% ==============================================================================
% SECTION 6: CROSS-ATTENTION WEIGHT EXTRACTION AND INTERPRETABILITY
% ==============================================================================
\section{Cross-attention weight extraction and interpretability}
\label{sec:lab02_cross_attention}

One of the most remarkable properties of attention-based architectures is their inherent \textbf{mechanistic interpretability}: by inspecting attention weight matrices, we can observe precisely which source tokens the model focuses upon when synthesizing each output token.

\subsection{Accessing internal cross-attention tensors}

To extract attention tensors in Hugging Face, we instantiate the model with \texttt{output\_attentions=True}:

\begin{lstlisting}[language=Python, caption={Extracting cross-attention weights during forward pass (Cells 90--96).}, label={lst:lab02_extract_attns}]
# Load model with attention extraction enabled
model = AutoModelForSeq2SeqLM.from_pretrained(model_name, output_attentions=True)

# Run forward pass for the first generation step
output = model(**tokens, decoder_input_ids=torch.tensor([[ tokenizer.pad_token_id ]]))

# output.cross_attentions is a tuple of length 12 (one tensor per decoder layer)
print("Number of cross-attention layers:", len(output.cross_attentions)) # 12

# Inspect tensor dimensions of Layer 0
layer0_attn = output.cross_attentions[0]
print("Cross-attention shape:", layer0_attn.shape)
# Output: torch.Size([1, 12, 1, 13])
\end{lstlisting}

The 4D shape $[B, H, T_{\text{dec}}, T_{\text{enc}}] = [1, 12, 1, 13]$ reflects:
\begin{itemize}
    \item $B = 1$: Batch size.
    \item $H = 12$: Number of distinct attention heads in the decoder layer.
    \item $T_{\text{dec}} = 1$: Single query token in the decoder (the initial pad token generating \texttt{"Hallo"}).
    \item $T_{\text{enc}} = 13$: Thirteen tokens in the input prompt (\texttt{translate}, \texttt{english}, \texttt{to}, \texttt{german}, \texttt{:}, \texttt{hello}, \texttt{,}, \texttt{how}, \texttt{are}, \texttt{you}, \texttt{?}, \texttt{</s>}).
\end{itemize}

\subsection{Early vs late layer attention patterns}

Cells~98--104 of \texttt{text-01-t5.ipynb} plot the attention distributions across different layers and across all generated output tokens, shown in Figure~\ref{fig:cross_attention_plots}.

\begin{figure}[htbp]
\centering
\begin{subfigure}[b]{0.32\textwidth}
    \centering
    \includegraphics[width=\textwidth]{figures/lab02/cross_attention_layer0.png}
    \caption{Layer 0 cross-attention.}
    \label{fig:cross_attn_l0}
\end{subfigure}
\hfill
\begin{subfigure}[b]{0.65\textwidth}
    \centering
    \includegraphics[width=\textwidth]{figures/lab02/cross_attention_all_layers.png}
    \caption{Cross-attention heatmaps across layers 0 to 9.}
    \label{fig:cross_attn_all}
\end{subfigure}
\vskip\baselineskip
\begin{subfigure}[b]{0.48\textwidth}
    \centering
    \includegraphics[width=\textwidth]{figures/lab02/cross_attention_full_sequence.png}
    \caption{Average cross-attention for full translation.}
    \label{fig:cross_attn_full}
\end{subfigure}
\caption{Cross-attention visualization in T5. Early layers attend broadly to task prefix instructions (\texttt{translate}, \texttt{english}, \texttt{german}), whereas deeper layers (6--11) dynamically align target tokens (\texttt{Hallo}) with their source counterparts (\texttt{hello}).}
\label{fig:cross_attention_plots}
\end{figure}

The empirical findings from Figure~\ref{fig:cross_attention_plots} reveal a distinct functional hierarchy across Transformer depth:
\begin{enumerate}
    \item \textbf{Shallow Layers (Layers 0--3)}: Attention heads focus heavily on the \textbf{task prefix prompt tokens} (\texttt{"translate"}, \texttt{"english"}, \texttt{"to"}, \texttt{"german"}). The model is determining the functional modality of the task (i.e., routing computational state toward German translation rather than English summarization).
    \item \textbf{Deep Layers (Layers 6--11)}: Attention heads dynamically shift focus to the \textbf{content tokens}. When generating \texttt{"Hallo"}, maximum attention weight concentrates sharply on \texttt{"hello"}. When generating \texttt{"wie"}, attention shifts to \texttt{"how"}.
\end{enumerate}

% ==============================================================================
% SECTION 7: SCALED DOT-PRODUCT ATTENTION FROM SCRATCH
% ==============================================================================
\section{Scaled dot-product attention from scratch: Mathematical foundations}
\label{sec:lab02_scaled_dot_attention}

Having explored Hugging Face's production implementation of T5, we transition to Module~2 of the laboratory (\texttt{text-02-attention.ipynb}): isolating and constructing the \textbf{Scaled Dot-Product Attention mechanism} from first principles in pure PyTorch.

\subsection{Query, Key, and Value vectors: Database retrieval analogy}

The attention mechanism models sequence interactions using the conceptual framework of a soft, differentiable \textbf{database query-retrieval system}~\cite{vaswani2017}:

\begin{intuition}[The Query, Key, Value Information Retrieval Analogy]
Consider searching for a video on YouTube:
\begin{itemize}
    \item \textbf{Query ($\vq$)}: The text string you type into the search bar. It represents what the current token is \textit{looking for}.
    \item \textbf{Key ($\vk$)}: The video title and metadata tags stored on YouTube's servers. It represents the \textit{indexing handle} describing what each candidate token contains.
    \item \textbf{Value ($\vv$)}: The actual video content streamed back to your screen. It represents the \textit{substantive information} propagated forward if a match occurs.
\end{itemize}
In self-attention, we evaluate the compatibility between Query $\vq_i$ and Key $\vk_j$ via their dot product $\vq_i^\top \vk_j$. Passing these compatibilities through a softmax yields normalized weights $\alpha_{ij}$. The output representation is the weighted mixture of Value vectors: $\vo_i = \sum_j \alpha_{ij} \vv_j$.
\end{intuition}

\subsection{Mathematical derivation and the scaling factor $\frac{1}{\sqrt{\dk}}$}

Given query matrix $\mQ \in \R^{T \times \dk}$, key matrix $\mK \in \R^{T \times \dk}$, and value matrix $\mV \in \R^{T \times \dv}$, the scaled dot-product attention equation is defined as:
\begin{equation}
    \Attn{\mQ, \mK, \mV} = \softmax\left(\frac{\mQ \mK^\top}{\sqrt{\dk}} + \mM\right) \mV
    \label{eq:scaled_dot_prod_attn}
\end{equation}

\begin{theorem}{Variance Preservation of the Scaling Factor}{variance_scaling}
Let components of query vector $\vq = [q_1, \dots, q_{\dk}]$ and key vector $\vk = [k_1, \dots, k_{\dk}]$ be independent, identically distributed random variables with zero mean and unit variance:
\begin{equation}
    \E[q_m] = \E[k_m] = 0, \quad \Var(q_m) = \Var(k_m) = 1, \quad \forall m \in \{1, \dots, \dk\}
\end{equation}
Then the dot product $S = \vq^\top \vk = \sum_{m=1}^{\dk} q_m k_m$ has expected value:
\begin{equation}
    \E[S] = \sum_{m=1}^{\dk} \E[q_m k_m] = \sum_{m=1}^{\dk} \E[q_m] \E[k_m] = 0
\end{equation}
and variance:
\begin{align}
    \Var(S) &= \sum_{m=1}^{\dk} \Var(q_m k_m) \nonumber \\
    &= \sum_{m=1}^{\dk} \left(\Var(q_m)\Var(k_m) + \Var(q_m)\E[k_m]^2 + \Var(k_m)\E[q_m]^2\right) \nonumber \\
    &= \sum_{m=1}^{\dk} (1 \cdot 1 + 0 + 0) = \sum_{m=1}^{\dk} 1 = \dk
\end{align}
Consequently, scaling the dot product by $\frac{1}{\sqrt{\dk}}$ yields a stabilized score $Z = \frac{\vq^\top \vk}{\sqrt{\dk}}$ with strictly unit variance:
\begin{equation}
    \Var(Z) = \Var\left(\frac{S}{\sqrt{\dk}}\right) = \frac{1}{\dk} \Var(S) = \frac{\dk}{\dk} = 1
\end{equation}
\end{theorem}

\begin{examinsight}[Softmax Saturation and Vanishing Gradients]
Why is dividing by $\sqrt{\dk}$ mandatory when $\dk$ is large (e.g., $\dk = 64$ or $128$)?

If dot products are not scaled, for large $\dk$ the magnitude of scores $|S|$ grows proportionally to $\sqrt{\dk}$. When large values (e.g., $+30$ or $-30$) enter the $\softmax$ function:
\begin{equation}
    \softmax(\vz)_i = \frac{e^{z_i}}{\sum_j e^{z_j}}
\end{equation}
the exponential function causes the probability distribution to collapse into a one-hot distribution (the maximum value receives probability $\approx 1$, all others $\approx 0$).

In this saturated region, the Jacobian derivative of the softmax function approaches zero:
\begin{equation}
    \frac{\partial \softmax(\vz)_i}{\partial z_j} = \softmax(\vz)_i (\delta_{ij} - \softmax(\vz)_j) \approx 0
\end{equation}
This induces severe \textbf{vanishing gradients}, preventing the optimizer from updating upstream query and key projection matrices during backpropagation! Dividing by $\sqrt{\dk}$ prevents saturation and stabilizes learning dynamics.
\end{examinsight}

\subsection{TikZ computation graph of Scaled Dot-Product Attention}

\begin{figure}[htbp]
\centering
\begin{tikzpicture}[
    scale=0.9,
    every node/.style={transform shape},
    tblock/.style={rectangle, draw=politoNavy, fill=skyBlue!30, very thick, rounded corners=4pt, minimum width=26mm, minimum height=9mm, font=\small\bfseries, align=center},
    opblock/.style={circle, draw=deepdiveTeal, fill=deepdiveTealBg, very thick, minimum size=10mm, font=\small\bfseries, align=center},
    mathblock/.style={rectangle, draw=crimsonRed, fill=examAmberBg, very thick, rounded corners=4pt, minimum width=28mm, minimum height=9mm, font=\small\bfseries, align=center},
    arr/.style={-{Stealth[scale=1.1]}, thick, draw=bodyDark}
]

\node[tblock] (Q) at (-3, 0) {Query $\mQ$\\$[B, T, \dk]$};
\node[tblock] (K) at (0, 0) {Key $\mK$\\$[B, T, \dk]$};
\node[tblock] (V) at (3.5, 0) {Value $\mV$\\$[B, T, \dv]$};

\node[opblock] (matmul1) at (-1.5, 1.8) {$\times$};
\node[font=\scriptsize\bfseries, above=0.1cm of matmul1] {$\mQ \mK^\top \in \R^{B \times T \times T}$};

\node[mathblock] (scale) at (-1.5, 3.2) {Scale by $\frac{1}{\sqrt{\dk}}$};
\node[mathblock] (mask) at (-1.5, 4.4) {Add Mask $\mM$\\+ Rel Bias $\mB$};
\node[mathblock] (softmax) at (-1.5, 5.6) {Softmax\\$\softmax(\cdot)$};

\node[opblock] (matmul2) at (1.0, 6.8) {$\times$};
\node[tblock] (out) at (1.0, 8.2) {Output $\mO$\\$[B, T, \dv]$};

\draw[arr] (Q) |- (matmul1);
\draw[arr] (K) |- (matmul1);
\draw[arr] (matmul1) -- (scale);
\draw[arr] (scale) -- (mask);
\draw[arr] (mask) -- (softmax);
\draw[arr] (softmax) |- (matmul2);
\draw[arr] (V) |- (matmul2);
\draw[arr] (matmul2) -- (out);

\end{tikzpicture}
\caption{Computational graph of the Scaled Dot-Product Attention mechanism with explicit tensor dimensions.}
\label{fig:scaled_dot_product_graph}
\end{figure}

Figure~\ref{fig:scaled_dot_product_graph} illustrates the explicit dataflow, showing how the intermediate score matrix $\mQ \mK^\top$ forms an all-pairs quadratic matrix of shape $[B, T, T]$ before being multiplied by $\mV$.

% ==============================================================================
% SECTION 8: SYNTHETIC TIME SERIES MODELING WITH PERIODIC JUMPS
% ==============================================================================
\section{Synthetic time series modeling with periodic jumps}
\label{sec:lab02_time_series}

To rigorously test the attention mechanism's ability to model sequential relationships, Module~2 formulates a challenging synthetic regression task: predicting the next value in a sequence exhibiting periodic, discontinuous step jumps.

\subsection{Problem formulation: Discontinuous periodic sequence prediction}

Standard autoregressive patterns (such as pure sinusoidal oscillations or linear trends) can often be solved trivially by linear autoregressive filters. To force the model to rely on non-local sequential memory, the dataset generates signals with abrupt jumps:
\begin{itemize}
    \item Sequence length: $T = 40$.
    \item A baseline signal initializes at zero.
    \item An initial jump of magnitude $5 + \epsilon$ occurs at a random offset $t_0 \in [0, 3)$.
    \item Subsequent jumps occur strictly every $\text{val} = 3$ timesteps ($t_0 + 3, t_0 + 6, t_0 + 9, \dots$), alternating between $0$ and $5 + \epsilon$.
\end{itemize}
The target $y \in \R$ is the signal value at timestep $T = 40$. To accurately predict $y$, the model cannot rely solely on the most recent step; it must determine whether timestep $40$ corresponds to a jump step relative to the sequence's periodic phase.

\subsection{Verified code: Custom \texttt{TimeSeriesDataset} and \texttt{DataLoader}}

Listing~\ref{lst:lab02_dataset_impl} displays the verified implementation from Cell~7 of \texttt{text-02-attention.ipynb}.

\begin{lstlisting}[language=Python, caption={Verified implementation of \texttt{TimeSeriesDataset} with periodic jumps (Cells 7--8).}, label={lst:lab02_dataset_impl}]
import numpy as np
import torch
import torch.nn as nn
import matplotlib.pyplot as plt

class TimeSeriesDataset(torch.utils.data.Dataset):
    def __init__(self, seq_len, num_samples):
        val = 3          # Periodicity of jumps
        initial_jump = 5 # Jump amplitude
        
        self.seq_len = seq_len
        self.data = []
        for _ in range(num_samples):
            time_series = np.zeros(seq_len + 1)
            
            # Introduce the first jump at a random point between 0 and val
            first_jump = np.random.randint(0, val)
            time_series[first_jump] += initial_jump + np.random.rand()
            
            # Continue introducing jumps every val timesteps after first jump
            for i in range(first_jump + val, seq_len + 1, val):
                jump = 0 if i % 2 == 1 else initial_jump
                time_series[i] += jump + np.random.rand()
            
            self.data.append(time_series)
    
    def __len__(self):
        return len(self.data)
    
    def __getitem__(self, idx):
        series = self.data[idx]
        # Return input sequence (first seq_len values) and target scalar (final value)
        return (torch.tensor(series[:self.seq_len], dtype=torch.float32), 
                torch.tensor(series[self.seq_len], dtype=torch.float32))

# Visualize a single generated sequence
dataset = TimeSeriesDataset(seq_len=40, num_samples=1)
x, y = dataset[0]
plt.figure(figsize=(7, 3))
plt.plot(range(40), x, label='x (input sequence)')
plt.scatter([40], [y], color='red', label=f'y (target: {y:.2f})')
plt.title(f"Synthetic Time Series with Periodic Jumps (Period=3)")
plt.legend()
plt.savefig("synthetic_time_series_sample.png")
\end{lstlisting}

\begin{figure}[htbp]
\centering
\includegraphics[width=0.72\textwidth]{figures/lab02/synthetic_time_series_sample.png}
\caption{Sample time series from the synthetic dataset generated in Laboratory~02. The sequence exhibits periodic spikes every 3 steps with randomized starting phase, requiring the attention mechanism to learn shift-invariant relative offsets.}
\label{fig:synthetic_sample}
\end{figure}

\subsection{Implementing custom \texttt{ScaledDotProductAttention} and \texttt{AttentionHead}}

Next, we construct the attention layers from scratch. Listing~\ref{lst:lab02_custom_attention} presents the verified implementations from Cells~14 and~18:

\begin{lstlisting}[language=Python, caption={Implementing \texttt{ScaledDotProductAttention} and \texttt{AttentionHead} (Cells 14--18).}, label={lst:lab02_custom_attention}]
class ScaledDotProductAttention(nn.Module):
    def __init__(self, seq_len, d_k, other_positional_encodings_present):
        super(ScaledDotProductAttention, self).__init__()
        self.d_k = d_k
        self.seq_len = seq_len
        self.other_positional_encodings_present = other_positional_encodings_present

        # Learnable relative positional encoding matrix [1, seq_len, seq_len]
        self.rel_layer = nn.Parameter(torch.randn(1, self.seq_len, self.seq_len))

    def forward(self, query, key, value):
        # 1. Compute dot products scaled by sqrt(d_k)
        scores = torch.matmul(query, key.transpose(-2, -1)) / math.sqrt(self.d_k)
        
        # 2. Add relative positional bias if absolute encodings are not present
        if not self.other_positional_encodings_present:
            scores = scores + self.rel_layer
        
        # 3. Softmax normalization across sequence dimension
        attention_weights = torch.softmax(scores, dim=-1)
        
        # 4. Multiply by Value matrix
        output = torch.matmul(attention_weights, value)
        return output, attention_weights


class AttentionHead(nn.Module):
    def __init__(self, d_model, seq_len, other_positional_encodings_present):
        super(AttentionHead, self).__init__()
        self.d_k = d_model
        self.seq_len = seq_len

        # Linear projections for Query, Key, and Value
        self.query_layer = nn.Linear(d_model, d_model)
        self.key_layer = nn.Linear(d_model, d_model)
        self.value_layer = nn.Linear(d_model, d_model)

        # Scaled dot-product core
        self.scaled_dot_attention = ScaledDotProductAttention(
            self.seq_len, self.d_k, other_positional_encodings_present
        )

        # Final projection layer
        self.fc_out = nn.Linear(d_model, d_model)

    def forward(self, query, key, value):
        Q = self.query_layer(query)
        K = self.key_layer(key)
        V = self.value_layer(value)

        attn_out, attn_weights = self.scaled_dot_attention(Q, K, V)
        output = self.fc_out(attn_out)
        return output, attn_weights
\end{lstlisting}

\subsection{Transformer Encoder Block and Time Series Architecture}

We assemble the attention head into a full \texttt{TransformerEncoderBlock} with residual connections, LayerNorm, and a feed-forward MLP (Listing~\ref{lst:lab02_encoder_block}), followed by the top-level \texttt{Transformer} regression model:

\begin{lstlisting}[language=Python, caption={Complete \texttt{TransformerEncoderBlock} and \texttt{Transformer} model (Cells 21--24).}, label={lst:lab02_encoder_block}]
class TransformerEncoderBlock(nn.Module):
    def __init__(self, d_model, dim_feedforward, seq_len, other_positional_encodings_present, dropout=0.1):
        super(TransformerEncoderBlock, self).__init__()
        self.attention = AttentionHead(d_model, seq_len, other_positional_encodings_present)
        self.norm1 = nn.LayerNorm(d_model)
        self.norm2 = nn.LayerNorm(d_model)

        # Feed-forward network
        self.ffn = nn.Sequential(
            nn.Linear(d_model, dim_feedforward),
            nn.ReLU(),
            nn.Linear(dim_feedforward, d_model)
        )
        self.dropout = nn.Dropout(dropout)

    def forward(self, x):
        # Multi-head attention with residual skip and LayerNorm
        attn_out, attn_weights = self.attention(x, x, x)
        x = self.norm1(x + self.dropout(attn_out))

        # Position-wise feed-forward with residual skip and LayerNorm
        ffn_out = self.ffn(x)
        output = self.norm2(x + self.dropout(ffn_out))
        return output, attn_weights


class Transformer(nn.Module):
    def __init__(self, input_size, d_model, dim_feedforward, num_layers, seq_len, 
                 positional_encoding=None, dropout=0.1, max_len=5000):
        super(Transformer, self).__init__()
        self.positional_encoding = positional_encoding

        # Linear projection: scalar time series feature -> d_model
        self.embedding = nn.Linear(input_size, d_model)
        
        if self.positional_encoding is not None:
            self.pos_encoder = self.positional_encoding(d_model, max_len)
        
        # Stacked Transformer Encoder blocks
        self.encoder_layers = nn.ModuleList([
            TransformerEncoderBlock(
                d_model, dim_feedforward, seq_len, 
                False if positional_encoding is None else True, 
                dropout
            ) for _ in range(num_layers)
        ])
        
        # Linear head projecting last timestep representation to predicted scalar
        self.fc_out = nn.Linear(d_model, 1)

    def forward(self, x):
        # x shape: [batch_size, seq_len, input_size]
        x = self.embedding(x)
        if self.positional_encoding is not None:
            x = self.pos_encoder(x)
        
        all_attention_weights = []
        for layer in self.encoder_layers:
            x, attn_weights = layer(x)
            all_attention_weights.append(attn_weights)

        # Predict from final timestep token (seq_len - 1)
        output = self.fc_out(x[:, -1])
        return output, all_attention_weights
\end{lstlisting}

% ==============================================================================
% SECTION 9: RELATIVE VS ABSOLUTE POSITIONAL ENCODINGS
% ==============================================================================
\section{Relative vs absolute positional encodings: Theory and empirical proof}
\label{sec:lab02_positional_encodings}

A central conceptual breakthrough of Laboratory~02 is understanding \textbf{Permutation Equivariance} and why the choice of positional representation dramatically dictates sequence learning capability.

\subsection{The permutation equivariance dilemma}

Standard self-attention is strictly \textbf{permutation equivariant}: if the rows of input matrix $\mX$ are permuted by permutation matrix $\mP$, the output is identically permuted:
\begin{equation}
    \Attn{\mP \mQ, \mP \mK, \mP \mV} = \mP \Attn{\mQ, \mK, \mV}
\end{equation}
In other words, without positional information, self-attention treats an input sequence as an unordered set (a bag of tokens)! To model sequential ordering, position must be injected into the architecture.

\subsection{Absolute sinusoidal positional encodings (Vaswani et al. 2017)}

In the original Transformer~\cite{vaswani2017}, Vaswani et al. added deterministic sinusoidal vectors to input embeddings:
\begin{align}
    PE_{(pos, 2i)} &= \sin\left(\frac{pos}{10000^{2i / \dmodel}}\right) \\
    PE_{(pos, 2i+1)} &= \cos\left(\frac{pos}{10000^{2i / \dmodel}}\right)
\end{align}
Listing~\ref{lst:lab02_sinusoidal_pe} displays the implementation from Cell~36:

\begin{lstlisting}[language=Python, caption={Absolute sinusoidal positional encoding implementation (Cell 36).}, label={lst:lab02_sinusoidal_pe}]
class PositionalEncoding(nn.Module):
    def __init__(self, d_model, max_len=5000):
        super(PositionalEncoding, self).__init__()
        pe = torch.zeros(max_len, d_model)
        position = torch.arange(0, max_len, dtype=torch.float).unsqueeze(1)
        div_term = torch.exp(torch.arange(0, d_model, 2).float() * (-math.log(10000.0) / d_model))
        
        pe[:, 0::2] = torch.sin(position * div_term)
        pe[:, 1::2] = torch.cos(position * div_term)
        
        pe = pe.unsqueeze(0).transpose(0, 1) # [max_len, 1, d_model]
        self.register_buffer('pe', pe)

    def forward(self, x):
        # Add absolute positional vector directly to embedded inputs
        return x + self.pe[:x.size(0), :].repeat(1, x.shape[1], 1)
\end{lstlisting}

\subsection{Relative positional encodings (Shaw et al. 2018 / T5)}

Rather than adding absolute coordinates to inputs, \textbf{Relative Positional Encodings}~\cite{shaw2018,raffel2020t5} recognize that what matters in language and sequence modeling is not whether a word is at index 37, but rather that word $j$ is \textbf{3 positions before} word $i$ ($i - j = 3$).

In our implemented relative attention layer, an offset bias matrix $\mB \in \R^{T \times T}$ is added directly to the attention logits:
\begin{equation}
    \text{Score}_{ij} = \frac{\vq_i^\top \vk_j}{\sqrt{\dk}} + B_{i, j}
\end{equation}
where $B_{i, j}$ directly models the pairwise relative distance between timesteps $i$ and $j$.

\subsection{Empirical convergence benchmark: Relative ($0.036$) vs Absolute ($5.417$)}

Both models were trained on identical synthetic datasets ($N = 1000$ sequences, sequence length $T = 40$, batch size $128$, Adam optimizer, learning rate $5 \times 10^{-4}$, MSE loss) across 15 epochs. 

The empirical loss trajectory recorded during Laboratory~02 is presented in Table~\ref{tab:relative_vs_absolute_loss} and plotted via PGFPlots in Figure~\ref{fig:loss_convergence_pgfplot}.

\begin{table}[htbp]
\centering
\small
\renewcommand{\arraystretch}{1.2}
\begin{tabular}{r c c}
\toprule
\textbf{Epoch} & \textbf{Relative Positional Encoding Loss} & \textbf{Absolute Sinusoidal Encoding Loss} \\
\midrule
1  & 7.3971 & 7.1401 \\
2  & 6.1423 & 6.5387 \\
3  & 5.3159 & 6.2578 \\
4  & 4.4852 & 5.9696 \\
5  & 3.7367 & 5.8090 \\
6  & 2.3744 & 5.7272 \\
7  & 1.0627 & 5.7824 \\
8  & 0.3977 & 5.6585 \\
9  & 0.2739 & 5.5258 \\
10 & 0.1367 & 5.6687 \\
11 & 0.0736 & 5.5961 \\
12 & 0.0533 & 5.4340 \\
13 & 0.0436 & 5.3413 \\
14 & 0.0380 & 5.3534 \\
\textbf{15} & \textbf{0.0363} & \textbf{5.4174} \\
\bottomrule
\end{tabular}
\caption{Empirical loss trajectory over 15 training epochs in Laboratory~02. The Relative Positional Encoding model achieves near-zero error ($0.0363$), while Absolute Sinusoidal Encodings completely fail to converge ($5.4174$).}
\label{tab:relative_vs_absolute_loss}
\end{table}

\begin{figure}[htbp]
\centering
\begin{tikzpicture}
\begin{axis}[
    width=0.88\textwidth,
    height=7.2cm,
    xlabel={\textbf{Training Epoch}},
    ylabel={\textbf{Mean Squared Error (MSE) Loss}},
    grid=major,
    grid style={dashed, lightGrayLine},
    xmin=1, xmax=15,
    ymin=0, ymax=8.0,
    legend pos=north east,
    legend cell align={left},
    every axis plot/.append style={thick}
]

% Relative Positional Encoding Loss Curve
\addplot[color=crimsonRed, mark=square*, mark size=2pt] coordinates {
    (1, 7.3971)
    (2, 6.1423)
    (3, 5.3159)
    (4, 4.4852)
    (5, 3.7367)
    (6, 2.3744)
    (7, 1.0627)
    (8, 0.3977)
    (9, 0.2739)
    (10, 0.1367)
    (11, 0.0736)
    (12, 0.0533)
    (13, 0.0436)
    (14, 0.0380)
    (15, 0.0363)
};
\addlegendentry{Relative Positional Encoding (Shaw/T5)}

% Absolute Positional Encoding Loss Curve
\addplot[color=steelBlue, mark=triangle*, mark size=2.5pt] coordinates {
    (1, 7.1401)
    (2, 6.5387)
    (3, 6.2578)
    (4, 5.9696)
    (5, 5.8090)
    (6, 5.7272)
    (7, 5.7824)
    (8, 5.6585)
    (9, 5.5258)
    (10, 5.6687)
    (11, 5.5961)
    (12, 5.4340)
    (13, 5.3413)
    (14, 5.3534)
    (15, 5.4174)
};
\addlegendentry{Absolute Sinusoidal Encoding (Vaswani)}

\end{axis}
\end{tikzpicture}
\caption{Analytical PGFPlots convergence benchmark from Laboratory~02. The model equipped with learnable relative bias converges smoothly from $7.3971 \to 0.0363$, successfully locking onto the period-3 invariant. The absolute sinusoidal model stagnates at $5.4174$ because randomized phase shifts destroy absolute index correlations.}
\label{fig:loss_convergence_pgfplot}
\end{figure}

\begin{figure}[htbp]
\centering
\begin{subfigure}[b]{0.48\textwidth}
    \centering
    \includegraphics[width=\textwidth]{figures/lab02/relative_attention_weights.png}
    \caption{Attention weights: Relative Positional Encoding.}
    \label{fig:rel_attn_weights}
\end{subfigure}
\hfill
\begin{subfigure}[b]{0.48\textwidth}
    \centering
    \includegraphics[width=\textwidth]{figures/lab02/absolute_attention_weights.png}
    \caption{Attention weights: Absolute Positional Encoding.}
    \label{fig:abs_attn_weights}
\end{subfigure}
\caption{Normalized attention weight profiles over input sequences. The relative model (a) concentrates sharp attention spikes exactly on periodic jump predecessors ($t - 3, t - 6$), whereas the absolute model (b) produces diffuse, uninformative attention spread.}
\label{fig:attn_weight_comparison}
\end{figure}

\subsection{Why relative attention dominates: Invariance under phase shifts}

The dramatic divergence in Figure~\ref{fig:loss_convergence_pgfplot} demonstrates a core principle of sequence modeling:
\begin{itemize}
    \item \textbf{Phase Invariance}: In our dataset, the initial jump occurs at a random offset $t_0 \in \{0, 1, 2\}$. Therefore, in Sample A, jumps might occur at indices $0, 3, 6, 9, \dots, 39$. In Sample B, jumps might occur at indices $2, 5, 8, 11, \dots, 38$.
    \item \textbf{Failure of Absolute Encodings}: The absolute model associates each step with fixed index coordinates ($PE_{pos}$). Because a jump occurs at index 39 in Sample A but at index 38 in Sample B, absolute index coordinates provide contradictory signals across samples. The model cannot generalize, causing training loss to plateau at $5.4174$.
    \item \textbf{Success of Relative Encodings}: The relative model evaluates pairwise distances $i - j$. In every single sample, regardless of where the sequence started, a jump at step $i$ is strictly correlated with what occurred at step $i - 3$! The relative bias matrix $B_{i, j}$ easily learns high affinity for relative offsets $\Delta = -3$ and $\Delta = -6$, allowing the model to achieve $0.0363$ loss.
\end{itemize}

\subsection{Modern extensions: RoPE and ALiBi}

While Shaw et al.~\cite{shaw2018} and T5 introduced relative position bias matrices, modern open-source foundation models deploy more advanced formulations:
\begin{enumerate}
    \item \textbf{Rotary Position Embeddings (RoPE)}~\cite{su2024roformer}: Deployed in LLaMA, Mistral, and Gemma. Instead of adding a bias to attention logits, RoPE rotates Query and Key vectors in the 2D complex plane by an angle proportional to their sequence position:
    \begin{equation}
        \vq_m = \mat{R}_{\Theta, m}^d \mW_q \vx_m, \quad \vk_n = \mat{R}_{\Theta, n}^d \mW_k \vx_n
    \end{equation}
    The inner product $\vq_m^\top \vk_n$ naturally depends strictly on the relative distance $m - n$, combining the simplicity of absolute input vectors with the theoretical advantages of relative position representations.
    \item \textbf{ALiBi (Attention with Linear Biases)}~\cite{press2022alibi}: Deployed in MPT and BLOOM. ALiBi adds a non-learned static negative penalty proportional to relative distance directly to attention logits:
    \begin{equation}
        \text{Score}_{ij} = \frac{\vq_i^\top \vk_j}{\sqrt{\dk}} - m \cdot |i - j|
    \end{equation}
    where $m$ is a head-specific slope. This achieves state-of-the-art sequence length extrapolation without requiring any positional parameter training.
\end{enumerate}

% ==============================================================================
% SECTION 10: INSTRUCTOR INSIGHTS, EXAM TAKEAWAYS, AND ENGINEERING CHECKLIST
% ==============================================================================
\section{Instructor insights, exam takeaways, and engineering checklist}
\label{sec:lab02_exam_insights}

\subsection{Key takeaways from the teaching assistant lecture walkthrough}

During the laboratory live session, teaching assistant Claudio Savelli emphasized several critical software engineering takeaways:
\begin{enumerate}
    \item \textbf{SentencePiece Detokenization}: The underscore prefix \texttt{\_} preserves whitespace boundary information. Decoding never guesses where spaces belonged; detokenization is lossless.
    \item \textbf{VRAM Saturation via Batch Outliers}: A single long sentence in a batch forces all other sequences to be zero-padded to that length, wasting substantial memory and GPU FLOPs. Always inspect sequence length histograms before batching.
    \item \textbf{Dimensional Matching in Multi-Head Layers}: Ensure that $\dmodel$ is cleanly divisible by $H$. In T5, $768 = 12 \times 64$.
    \item \textbf{Weight Sharing Verification}: T5 saves tens of millions of parameters by sharing a single tensor between encoder embeddings, decoder embeddings, and the language modeling classification head.
\end{enumerate}

\subsection{Recurrent exam questions and traps}

\begin{examinsight}[Common Exam Traps on Attention and Seq2Seq Models]
\begin{itemize}
    \item \textbf{Trap 1: Cross-Attention Dimensions}: On exams, students often incorrectly assume cross-attention has shape $[B, H, T_{\text{enc}}, T_{\text{enc}}]$. Remember: in cross-attention, Queries originate from the \textbf{decoder} ($T_{\text{dec}}$), while Keys and Values originate from the \textbf{encoder} ($T_{\text{enc}}$). The resulting attention matrix has shape:
    \begin{equation}
        [B, H, T_{\text{dec}}, T_{\text{enc}}]
    \end{equation}
    \item \textbf{Trap 2: Forgetting the Causal Mask in Autoregression}: If a student implements a custom decoder without applying lower-triangular masking ($M_{ij} = -\infty$ for $j > i$), the decoder cheats during training by looking ahead at ground truth future tokens, resulting in perfect training loss but completely corrupted inference generation.
    \item \textbf{Trap 3: Why Divide by $\sqrt{\dk}$}: Be prepared to write down Theorem~\ref{thm:variance_scaling}. You must demonstrate that without scaling, the dot-product variance scales with $\dk$, causing softmax saturation and vanishing backpropagation gradients.
    \item \textbf{Trap 4: Relative vs Absolute Positional Encodings}: Absolute encodings assign fixed coordinates to sequence indices; relative encodings model pairwise distance offsets $i - j$. In periodic or translation-invariant sequences with varying start offsets, absolute encodings fail while relative encodings succeed.
\end{itemize}
\end{examinsight}

\subsection{Engineering checklist: Validating Transformer implementations}

When building or fine-tuning Transformer models in PyTorch, verify each step against this engineering checklist:
\begin{itemize}
    \item[$\square$] \textbf{Padding Mask Integration}: Have you verified that pad tokens (\texttt{id=0}) have attention mask set to \texttt{0} in both encoder and decoder?
    \item[$\square$] \textbf{Scaling Factor Check}: Is your dot-product division strictly $\frac{1}{\sqrt{\dk}}$ and not $\frac{1}{\sqrt{\dmodel}}$?
    \item[$\square$] \textbf{KV-Cache Deployment}: For production inference loops, is \texttt{use\_cache=True} activated to avoid $O(T^2)$ recomputation?
    \item[$\square$] \textbf{Weight Tying Registration}: If tying embedding weights to classification heads, ensure they share identical memory addresses (\texttt{lm\_head.weight = shared.weight}).
    \item[$\square$] \textbf{Permutation Invariance Test}: If building custom attention architectures, pass a permuted sequence to verify whether positional encodings correctly break permutation symmetry.
\end{itemize}

\subsection{Bridge to Laboratory 03}
In Laboratory~03, we extend these foundational Transformer mechanics to full-scale modern language modeling architectures: delving deep into the bidirectional masked language modeling of BERT~\cite{devlin2019}, the causal autoregressive pretraining of GPT-2~\cite{radford2019}, and parameter-efficient fine-tuning on downstream software engineering tasks.

\end{document}
